\customchapter{RESUMEN}

La calidad del aire en espacios cerrados, especialmente en aulas universitarias, es crucial para el bienestar y el rendimiento de los estudiantes. Altos niveles de dióxido de carbono (CO\(_2\)) y material particulado (polvo), sumados a variaciones de temperatura y humedad, pueden afectar negativamente la salud y el desempeño cognitivo de los ocupantes. En este contexto, este trabajo presenta el diseño e implementación de un sistema de monitoreo ambiental multivariable y alerta en tiempo real basado en una arquitectura de red de Internet de las Cosas (IoT). El sistema está compuesto por nodos sensores que recopilan mediciones de CO\(_2\), polvo, temperatura y humedad, transmitiéndolas vía comunicación inalámbrica LoRa hacia un gateway personalizado basado en ESP32. Este gateway actúa como puente para enviar los datos hacia una infraestructura en la nube, desplegada en Microsoft Azure mediante Infraestructura como Código (IaC). La plataforma web para la visualización ha sido desarrollada utilizando React, acoplada a una API alojada en Azure Container Apps, y los datos se almacenan en Azure Cosmos DB. El sistema incluye la generación de alertas automatizadas y un subsistema de mitigación activa que compara el rendimiento de diversas estrategias de control de ventilación (tales como PID, entre otras) cuando se exceden los umbrales permisibles, contribuyendo a la creación de ambientes educativos más seguros. Las pruebas de integración validaron el flujo de datos completo y su capacidad para monitorear y accionar medidas de mitigación en tiempo real.

\noindent \textbf{Palabras clave:}\\
\noindent Calidad del aire; CO\(_2\); material particulado; temperatura; humedad; IoT; LoRa; Azure; control comparativo; ventilación automatizada.